\begin{corollary}[3.2.7]
\label{IV.2.3.2.7-fr}
Bajo las hipótesis de \textup{(3.2.6)}, si $\mathcal F$ carece de ciclo primo asociado
inmerso, solo existe una descomposición irredundante reducida de $\mathcal F$.
\end{corollary}

\begin{corollary}[3.2.8]
\label{IV.2.3.2.8-fr}
Sean $X$ un preesquema noetheriano y $\mathcal F$ un $\mathcal O_X$-módulo
coherente. Existe una filtración finita $(\mathcal F_i)_{0\le i\le n}$ de
$\mathcal F$ tal que $\mathcal F_0=\mathcal F$, $\mathcal F_n=0$, formada por
$\mathcal O_X$-módulos coherentes y tal que los cocientes
$\mathcal F_i/\mathcal F_{i+1}$ sean nulos o irredundantes, y
$\operatorname{Ass}(\mathcal F_i/\mathcal F_{i+1})
\subset\operatorname{Ass}(\mathcal F)$.
\end{corollary}

En efecto, $\mathcal F$ es isomorfo a un sub-$\mathcal O_X$-módulo de una
suma directa finita $\bigoplus_{j=1}^n\mathcal G_j$, donde los $\mathcal G_j$
son irredundantes y coherentes (3.2.6); como todo
sub-$\mathcal O_X$-módulo cuasicoherente de $\mathcal G_j$ es nulo o
irredundante (3.1.7), los
$\mathcal F_i=\mathcal F\cap(\bigoplus_{j=1}^{n-i}\mathcal G_j)$ satisfacen lo
pedido, siendo $\mathcal F_i/\mathcal F_{i+1}$ isomorfo a un
sub-$\mathcal O_X$-módulo coherente de $\mathcal G_{n-i}$.

